% lecture.tex
\documentclass[aspectratio=169, xcolor={table,xcdraw}]{beamer}

\usetheme{Madrid}
\usecolortheme{default}
\definecolor{ucbblue}{RGB}{0, 51, 102}

\setbeamercolor{structure}{fg=ucbblue}
\setbeamercolor*{palette primary}{fg=white,bg=ucbblue}
\setbeamercolor*{palette secondary}{fg=white,bg=ucbblue!80}
\setbeamercolor*{palette tertiary}{fg=white,bg=ucbblue!90}
\setbeamercolor{title}{fg=white, bg=ucbblue}
\setbeamercolor{frametitle}{fg=white, bg=ucbblue}
\setbeamercolor{item}{fg=ucbblue}
\setbeamercolor{block title}{bg=ucbblue!15, fg=ucbblue}
\setbeamercolor{block body}{bg=ucbblue!5}

\usepackage[T1]{fontenc}
\usepackage[utf8]{inputenc}
\usepackage{tgpagella}
\usefonttheme{serif}
\usepackage[spanish, es-noshorthands]{babel}
\usepackage{amsmath, amssymb, bm}
\usepackage{graphicx}
\usepackage[most]{tcolorbox}
\usepackage{booktabs}
\usepackage{tabularx}
\usepackage{tikz}
\usepackage{pgfplots}
\pgfplotsset{compat=1.18}
\usepackage[american]{circuitikz}
\usetikzlibrary{shapes.geometric, arrows.meta, positioning, calc, backgrounds}

\title[Fundamentos Dinámicos RC]{Fundamentos Dinámicos a través de una Red de Primer Orden}
\subtitle{IMT-221 Circuitos Electrónicos II — Sesión 1/10 (Tutoría Intensiva)}
\author[B. Quiroga Turdera]{M.Sc. Ing. Bernardo Quiroga Turdera}
\institute[UCB Tarija]{Universidad Católica Boliviana ``San Pablo'' — Sede Tarija}
\date{Sesión 1}

\begin{document}

\begin{frame}
    \titlepage
\end{frame}

\begin{frame}{Arquitectura del Proyecto a Futuro}
    \begin{tcolorbox}[colback=ucbblue!5, colframe=ucbblue, title=V1: ¿Hacia dónde vamos?]
        \footnotesize Hoy caracterizamos un circuito RC autónomo. Este banco es un paso metodológico para predecir, simular y medir dinámicas; nos prepara para el \textbf{sistema de regulación térmica}.
    \end{tcolorbox}
    
    \begin{center}
    \begin{tikzpicture}[node distance=0.5cm and 0.4cm, thick,
        block/.style={draw, rounded corners, fill=blue!10, text width=1.6cm, align=center, minimum height=0.9cm, font=\scriptsize},
        highlight/.style={draw=red, line width=1.2pt, fill=red!10, text width=1.8cm, align=center, minimum height=0.9cm, font=\scriptsize}]
        
        % Fila superior (Camino directo)
        \node (consigna) [font=\scriptsize] {Consigna};
        \node[draw, circle, right=0.2cm of consigna, inner sep=1pt] (sum) {\tiny $\Sigma$};
        \node[block, right=0.3cm of sum] (control) {PI/PID};
        \node[block, right=0.3cm of control] (potencia) {Interfaz\\Potencia};
        \node[block, right=0.3cm of potencia] (planta) {Planta\\Térmica};
        
        % Fila inferior (Lazo de realimentación)
        \node[block, below=0.5cm of planta] (sensor) {Sensor};
        \node[block, left=0.3cm of sensor] (acond) {Acondic.};
        \node[highlight, left=0.3cm of acond] (filtro) {Filtro RC\\(Sesión 1)};
        
        % Conexiones
        \draw[->] (consigna) -- node[above, font=\tiny]{$+$} (sum);
        \draw[->] (sum) -- (control);
        \draw[->] (control) -- (potencia);
        \draw[->] (potencia) -- (planta);
        \draw[->] (planta) -- (sensor);
        \draw[->] (sensor) -- (acond);
        \draw[->] (acond) -- (filtro);
        \draw[->] (filtro) -| node[pos=0.95, left, font=\tiny]{$-$} (sum);
        
        % Perturbación
        \draw[->, dashed] ([yshift=0.4cm]planta.north) -- node[left, font=\tiny]{Perturb.} (planta.north);
    \end{tikzpicture}
    \end{center}
\end{frame}

\begin{frame}{Representación de Sinusoides y Fasores}
    \begin{columns}[T]
        \begin{column}{0.48\textwidth}
            Una señal sinusoidal física se describe por:
            \begin{equation*}
                v(t) = \hat{V} \cos(\omega t + \theta)
            \end{equation*}
            \begin{itemize}
                \item $\hat{V}$: Amplitud de pico
                \item $\omega = 2\pi f$: Frecuencia angular [rad/s]
                \item $\theta$: Fase inicial
                \item $T = 1/f$: Período
            \end{itemize}
            El \textbf{fasor} codifica amplitud y fase:
            \begin{equation*}
                \tilde{V} = \hat{V} e^{j\theta} \implies v(t) = \text{Re}\{\tilde{V} e^{j\omega t}\}
            \end{equation*}
        \end{column}
        \begin{column}{0.48\textwidth}
            \centering
            \textbf{V2: Medición Temporal de Fase}
            \vspace{0.2cm}
            
            \begin{tikzpicture}
                \begin{axis}[
                    width=6.5cm, height=4.2cm,
                    axis lines=middle,
                    xmin=0, xmax=10, ymin=-1.5, ymax=1.5,
                    xlabel={$t$}, ylabel={$v(t)$},
                    xtick=\empty, ytick={1}, yticklabels={$\hat{V}$},
                    domain=0:10, samples=100
                ]
                \addplot[thick, blue] {sin(deg(x))} node[pos=0.8, above] {$v_{in}$};
                \addplot[thick, red] {0.8*sin(deg(x - 1))} node[pos=0.9, right] {$v_{out}$};
                
                \draw[dashed, gray] (axis cs: 3.1415, 0) -- (axis cs: 3.1415, -0.5);
                \draw[dashed, gray] (axis cs: 4.1415, 0) -- (axis cs: 4.1415, -0.5);
                \draw[<->, thick] (axis cs: 3.1415, -0.4) -- node[below]{\small $\Delta t$} (axis cs: 4.1415, -0.4);
                \node[below] at (axis cs: 7, -0.8) {\small $\phi = -360^\circ f \Delta t$};
                \end{axis}
            \end{tikzpicture}
        \end{column}
    \end{columns}
\end{frame}

\begin{frame}{Impedancia del Capacitor}
    \begin{columns}[T]
        \begin{column}{0.54\textwidth}
            La relación voltaje-corriente en el tiempo es: \par
            {\centering $\displaystyle i_C(t) = C \frac{dv_C(t)}{dt}$ \par}
            \vspace{0.2cm}
            Sustituyendo la excitación fasorial: \par
            {\centering $\displaystyle \tilde{I}_C = j\omega C \tilde{V}_C$ \par}
            \vspace{0.2cm}
            La \textbf{impedancia capacitiva} es el cociente: \par
            {\centering $\displaystyle Z_C(j\omega) = \frac{\tilde{V}_C}{\tilde{I}_C} = \frac{1}{j\omega C} = -j\frac{1}{\omega C}$ \par}
            \vspace{0.2cm}
            \begin{itemize}
                \item Magnitud: $\vert{}Z_C\vert{} = \frac{1}{\omega C} = \frac{1}{2\pi f C}$
                \item Fase: $\arg(Z_C) = -90^\circ$
            \end{itemize}
        \end{column}
        \begin{column}{0.44\textwidth}
            \centering
            \textbf{V4: $\vert{}Z_C\vert{}$ contra Frecuencia}
            \vspace{0.2cm}
            
            \begin{tikzpicture}
                \begin{loglogaxis}[
                    width=6cm, height=4.5cm,
                    xlabel={$f$ (Hz)}, ylabel={$|Z_C|$ ($\Omega$)},
                    grid=both,
                    xmin=10, xmax=1000, ymin=1000, ymax=100000,
                    domain=10:1000
                ]
                \addplot[thick, blue, samples=50] {1/(2*pi*x*100e-9)};
                \addplot[mark=*, red] coordinates {(100, 15915)} node[above right] {$f_1$};
                \addplot[mark=*, red] coordinates {(1000, 1591)} node[above right] {$f_2$};
                \end{loglogaxis}
            \end{tikzpicture}
        \end{column}
    \end{columns}
\end{frame}

\begin{frame}{Derivación Analítica del Filtro Pasa-Bajos}
    \begin{columns}[T]
        \begin{column}{0.38\textwidth}
            \centering
            \textbf{V3: Topología del Filtro}
            \vspace{0.3cm}
            
            \begin{circuitikz}[thick, scale=0.75, transform shape]
                \draw (0,2) to[short, o-] (1,2) to[R, l=$R$] (3,2) to[short, -o] (4.5,2);
                \draw (3,2) to[C, l=$C$] (3,0);
                \draw (0,0) to[short, o-o] (4.5,0);
                \draw (3,0) node[ground]{};
                
                \node[left] at (0,2) {$v_{in}(t)$};
                \node[right] at (4.5,2) {$v_{out}(t)$};
            \end{circuitikz}
        \end{column}
        \begin{column}{0.60\textwidth}
            \small
            \textbf{1. Ley de Corrientes (KCL):} \par
            Corriente en el resistor iguala corriente en el capacitor:
            {\centering $\displaystyle \frac{v_{in}(t) - v_{out}(t)}{R} = C \frac{dv_{out}(t)}{dt}$ \par}
            \vspace{0.2cm}
            
            Reorganizando y definiendo constante $\tau = RC$:
            {\centering $\displaystyle \tau \frac{dv_{out}(t)}{dt} + v_{out}(t) = v_{in}(t)$ \par}
            \vspace{0.4cm}
            
            \textbf{2. Divisor de Impedancias ($s = j\omega$):}
            {\centering $\displaystyle H(s) = \frac{V_{out}(s)}{V_{in}(s)} = \frac{1/(sC)}{R + 1/(sC)}$ \par}
            \vspace{0.2cm}
            
            Multiplicando numerador y denominador por $sC$:
            {\centering $\displaystyle H(s) = \frac{1}{1 + sRC} = \frac{1}{1 + s\tau}$ \par}
        \end{column}
    \end{columns}
\end{frame}

\begin{frame}{Respuesta en Frecuencia: Normalización Analítica}
    Evaluamos la función de transferencia en régimen permanente senoidal ($s = j\omega$):
    {\centering $\displaystyle H(j\omega) = \frac{1}{1 + j\omega RC}$ \par}
    \vspace{0.3cm}
    
    Definimos la \textbf{frecuencia de corte} como el punto de diseño donde $\omega_c RC = 1$:
    {\centering $\displaystyle \omega_c = \frac{1}{RC} \text{ [rad/s]} \quad \implies \quad f_c = \frac{1}{2\pi RC} \text{ [Hz]}$ \par}
    \vspace{0.3cm}
    
    Normalizando la frecuencia física respecto a la frecuencia de corte ($x = f/f_c = \omega RC$), obtenemos la \textbf{forma canónica} del sistema de primer orden:
    {\centering $\displaystyle H(j\omega) = \frac{1}{1 + jx} \cdot \frac{1 - jx}{1 - jx} = \frac{1 - jx}{1 + x^2}$ \par}
\end{frame}

\begin{frame}{Ecuaciones Analíticas de Bode (Magnitud y Fase)}
    Extrayendo el módulo y el argumento de la forma canónica compleja, obtenemos las ecuaciones de predicción para el diagrama de Bode ideal:
    
    \begin{block}{Comportamiento en Régimen Permanente}
        \begin{align*}
            \text{Ganancia de tensión:} \quad & G(f) = |H(j\omega)| = \frac{1}{\sqrt{1 + (f/f_c)^2}} \\[2ex]
            \text{Ganancia en Decibelios:} \quad & G_{dB}(f) = 20 \log_{10} G(f) = -10 \log_{10}\left[1 + \left(\frac{f}{f_c}\right)^2\right] \\[2ex]
            \text{Desfase (Salida vs Entrada):} \quad & \phi(f) = \arg\{H(j\omega)\} = -\arctan\left(\frac{f}{f_c}\right)
        \end{align*}
    \end{block}
\end{frame}

\begin{frame}{V5: Diagrama de Bode Analítico Ideal}
    \vspace*{-0.2cm}
    \centering
    \begin{tikzpicture}
        \begin{semilogxaxis}[
            name=mag,
            width=0.85\textwidth, height=1.8cm,
            scale only axis,
            ylabel={$G_{dB}$ (dB)},
            xticklabel=\empty,
            ymin=-25, ymax=5,
            ytick={-20, -3, 0},
            grid=both,
            domain=0.01:100,
            samples=100,
            label style={font=\small},
            tick label style={font=\footnotesize}
        ]
            \addplot[thick, blue] {-10*log10(1 + x^2)};
            \addplot[dashed, red] coordinates {(0.01, 0) (1, 0)};
            \addplot[dashed, red] { -20*log10(x) };
            \addplot[mark=*, black, mark size=1.5pt] coordinates {(1, -3.01)} node[below left, font=\scriptsize] {$-3\text{ dB}$};
        \end{semilogxaxis}
        
        \begin{semilogxaxis}[
            at={(mag.south west)},
            anchor=north west,
            yshift=-0.5cm,
            width=0.85\textwidth, height=1.8cm,
            scale only axis,
            xlabel={$f/f_c$ (Frecuencia Normalizada)},
            ylabel={$\phi$ ($^\circ$)},
            ymin=-95, ymax=5,
            ytick={0, -45, -90},
            grid=both,
            domain=0.01:100,
            samples=100,
            label style={font=\small},
            tick label style={font=\footnotesize}
        ]
            \addplot[thick, blue] {-atan(x)};
            \addplot[mark=*, black, mark size=1.5pt] coordinates {(1, -45)} node[above right, font=\scriptsize] {$-45^\circ$};
        \end{semilogxaxis}
    \end{tikzpicture}
\end{frame}

\begin{frame}[fragile]{Metodología de Validación: Simulación y Laboratorio}
    \begin{columns}[T]
        \begin{column}{0.4\textwidth}
            \textbf{Starter SPICE (.ac):}
\begin{verbatim}
* CP1A - RC low pass 
V1 in 0 AC 1
R1 in out 10k
C1 out 0 100n
.ac dec 100 10 10000
.end
\end{verbatim}
            \vspace{0.2cm}
            \small \textit{Graficar $V(out)/V(in)$ para observar la transferencia $H(s)$ independientemente de la fuente física.}
        \end{column}
        \begin{column}{0.58\textwidth}
            \centering
            \textbf{V6: Flujo de Contraste y Discrepancia}\\
            \vspace{0.2cm}
            \begin{tikzpicture}[node distance=0.8cm, thick,
                block/.style={draw, rounded corners, fill=blue!10, text width=2cm, align=center, minimum height=1.1cm, font=\scriptsize}]
                
                \node[block] (pred) {\textbf{Predicción}\\Ecuación};
                \node[block, right=0.3cm of pred] (sim) {\textbf{Simulación}\\Barrido AC};
                \node[block, right=0.3cm of sim] (med) {\textbf{Medición}\\Física};
                \node[block, draw=red, fill=red!10, below=0.4cm of sim] (exp) {\textbf{Hipótesis}\\Explicación};
                
                \draw[->] (pred) -- (sim);
                \draw[->] (sim) -- (med);
                \draw[->] (med) |- (exp);
                \draw[->] (pred) |- (exp);
            \end{tikzpicture}
            
            \vspace{0.3cm}
            \begin{tcolorbox}[colback=white, colframe=gray, width=\textwidth, center]
                \centering \footnotesize
                \textbf{Cuantificación del Error (Misma Frecuencia):} \\
                \vspace{0.1cm}
                $E_G(\%) = 100 \frac{\vert G_{meas} - G_{pred}\vert}{\vert G_{pred}\vert} \quad \Delta \phi = \phi_{meas} - \phi_{pred}$
            \end{tcolorbox}
        \end{column}
    \end{columns}
\end{frame}

\begin{frame}{V8: Analogía de Redes (Eléctrico vs. Térmico)}
    \textbf{Advertencia Científica:} El modelo térmico es un \textbf{candidato} pendiente de validación en la Sesión 2.
    
    \vspace{0.2cm}
    \begin{columns}[T]
        \begin{column}{0.48\textwidth}
            \centering
            \textbf{Sistema Eléctrico RC} \\
            (Validado en Sesión 1)
            \begin{equation*}
                \tau_e \frac{dv_{out}}{dt} + v_{out} = v_{in}
            \end{equation*}
            \begin{equation*}
                H_e(s) = \frac{1}{1 + s\tau_e}
            \end{equation*}
            \begin{itemize}
                \item Ganancia DC: $1$ (V/V) \textit{adimensional}.
                \item Dinámica: $\tau_e = RC$ [s].
            \end{itemize}
        \end{column}
        \begin{column}{0.48\textwidth}
            \centering
            \textbf{Sistema Térmico Candidato} \\
            (Planta Proyecto)
            \begin{equation*}
                C_{th} \frac{d\Delta T}{dt} + G_{th}\Delta T = P(t)
            \end{equation*}
            \begin{equation*}
                \frac{\Delta T(s)}{P(s)} = \frac{1/G_{th}}{1 + s \frac{C_{th}}{G_{th}}}
            \end{equation*}
            \begin{itemize}
                \item Ganancia DC: $1/G_{th}$ [K/W].
                \item Dinámica: $\tau_{th} = C_{th}/G_{th}$ [s].
            \end{itemize}
        \end{column}
    \end{columns}
    \vspace{0.2cm}
    \centering
    \textit{Misma estructura matemática de 1\textsuperscript{er} orden diferencial, diferente física.}
\end{frame}

\end{document}
